\section{Conclusion}

In this paper, the sensitivity of two-agent EF$^{21}$ to heterogeneous local regularity is studied using HAC Law reported in Ref.~\cite{thomsen2026tight}. Before the sensitivity analysis, the evidential basis of the relation is examined using the source PEP and Lyapunov verification protocol and an independent reproduction of the complete 10,500-configuration cubic grid. The replay reaches a maximum polynomial residual of $1.742\times10^{-15}$ and recovers the homogeneous theoretical limit within $8.389\times10^{-13}$. These checks support numerical use of the cubic on the audited domain while preserving its empirical status.

The equal-smoothness study shows that strong-convexity heterogeneity reduces the available contraction margin and that compression increases its relative effect. The full-regularity study gives a different view of the same problem. With the arithmetic means of smoothness and strong convexity fixed, the empirical step size and $K_2$ remain unchanged, while the remaining coefficient gap is the weighted dispersion of the local condition-shape coordinates. This gap contains the squared mismatch $(\tau_L-\tau_\mu)^2$.

When the two heterogeneity ratios are aligned, the workers may still have different local regularity scales, but the cubic is unchanged from the homogeneous controlled case. A nonzero mismatch makes the gap positive and increases the predicted contraction factor. The generic root analysis also shows that the selected largest root increases strictly with $K_1$ at fixed $K_2$. The results therefore distinguish the amount of local heterogeneity from the mismatch between the two regularity ratios.

Further work may examine whether a similar mismatch coordinate appears for $n>2$, either through performance-estimation methods or a separate convergence analysis. Stochastic large-scale experiments and the possible use of the mismatch quantity in adaptive weighting or compression-aware design also remain open directions.

\authorcontributions{Author contributions will be finalized using the CRediT taxonomy before submission. All authors have read and agreed to the published version of the manuscript.}

\funding{Funding information will be confirmed before submission.}

\institutionalreview{Not applicable. This study did not involve human participants or animals.}

\informedconsent{Not applicable.}

\dataavailability{
The analysis code, numerical summaries, symbolic-verification materials,
and reproducibility scripts supporting this study are available in the
public GitHub repository at
\href{https://github.com/Jaycee871/Sensitivity-of-TwoAgent-EF21-to-Heterogeneous-Regularity-under-Compression}
{Sensitivity of Two-Agent EF21 to Heterogeneous Regularity under Compression}.
}

\acknowledgments{During the preparation of this study and manuscript, the authors used ChatGPT (OpenAI; GPT-5.6 Sol) to assist with research-workflow design, code generation and review, computational experiment planning, interpretation checks, and manuscript drafting and editing. Wolfram Language 15.0.1 was used as an independent symbolic-computation environment for algebraic simplification and discriminant analysis. All AI-assisted outputs, code, numerical and symbolic results, citations, and manuscript text were reviewed and verified by the human authors, who take full responsibility for the content of the publication.}

\conflictsofinterest{Fu-Hsing Wang serves as a Guest Editor of the Special Issue ``Artificial Intelligence and Algorithms'' in \emph{Mathematics}. Independent editorial handling by an Editorial Board member without a conflict of interest is requested. Any additional conflicts of interest will be confirmed and declared before submission.}

\abbreviations{Abbreviations}{The following abbreviations are used in this manuscript:\\
\noindent\begin{tabular}{@{}ll}
EF & Error Feedback\\
EF21 & Error Feedback 21\\
PEP & Performance Estimation Problem\\
\end{tabular}}

\reftitle{References}
\begin{thebibliography}{99}

\bibitem{seide2014onebit}
Seide, F.; Fu, H.; Droppo, J.; Li, G.; Yu, D.
1-bit stochastic gradient descent and its application to data-parallel distributed training of speech DNNs.
In \emph{Proceedings of Interspeech 2014};
Singapore, 14--18 September 2014;
pp. 1058--1062.
\href{https://doi.org/10.21437/Interspeech.2014-274}
{doi:10.21437/Interspeech.2014-274}.

\bibitem{stich2018sparsified}
Stich, S.U.; Cordonnier, J.-B.; Jaggi, M.
Sparsified SGD with memory.
In \emph{Advances in Neural Information Processing Systems};
2018; Volume 31, pp. 4452--4463.

\bibitem{karimireddy2019error}
Karimireddy, S.P.; Rebjock, Q.; Stich, S.U.; Jaggi, M.
Error feedback fixes SignSGD and other gradient compression schemes.
In \emph{Proceedings of the 36th International Conference on Machine Learning};
PMLR, 2019; Volume 97, pp. 3252--3261.

\bibitem{richtarik2021ef21}
Richt\'arik, P.; Sokolov, I.; Fatkhullin, I.
EF21: A new, simpler, theoretically better, and practically faster error feedback.
In \emph{Advances in Neural Information Processing Systems};
2021; Volume 34, pp. 4384--4396.

\bibitem{thomsen2026tight}
Berg Thomsen, D.; Taylor, A.B.; Dieuleveut, A.
A Tight Theory of Error Feedback Algorithms in Distributed Optimization.
\emph{arXiv} \textbf{2026}, arXiv:2605.31594.
\href{https://doi.org/10.48550/arXiv.2605.31594}
{\nolinkurl{doi:10.48550/arXiv.2605.31594}}.

\bibitem{drori2014performance}
Drori, Y.; Teboulle, M.
Performance of first-order methods for smooth convex minimization: A novel approach.
\emph{Math. Program.} \textbf{2014}, \emph{145}, 451--482.
\href{https://doi.org/10.1007/s10107-013-0653-0}
{doi:10.1007/s10107-013-0653-0}.

\bibitem{taylor2017exact}
Taylor, A.B.; Hendrickx, J.M.; Glineur, F.
Exact worst-case performance of first-order methods for composite convex optimization.
\emph{SIAM J. Optim.} \textbf{2017}, \emph{27}, 1283--1313.
\href{https://doi.org/10.1137/16M108104X}
{doi:10.1137/16M108104X}.

\bibitem{goujaud2024pepit}
Goujaud, B.; Moucer, C.; Glineur, F.; Hendrickx, J.M.; Taylor, A.B.; Dieuleveut, A.
PEPit: Computer-assisted worst-case analyses of first-order optimization methods in Python.
\emph{Math. Program. Comput.} \textbf{2024}, \emph{16}, 337--367.
\href{https://doi.org/10.1007/s12532-024-00259-7}
{doi:10.1007/s12532-024-00259-7}.

\bibitem{richtarik2024reloaded}
Richt\'arik, P.; Gasanov, E.; Burlachenko, K.
Error Feedback Reloaded: From Quadratic to Arithmetic Mean of Smoothness Constants.
In \emph{Proceedings of the 12th International Conference on Learning Representations (ICLR 2024)};
2024.

\bibitem{gao2024econtrol}
Gao, Y.; Islamov, R.; Stich, S.U.
EControl: Fast Distributed Optimization with Compression and Error Control.
In \emph{Proceedings of the 12th International Conference on Learning Representations (ICLR 2024)};
2024.

\bibitem{colla2026symmetries}
Colla, S.; Hendrickx, J.M.
Exploiting Agent Symmetries for Performance Analysis of Distributed Optimization Methods.
\emph{Open J. Math. Optim.} \textbf{2026}, \emph{7}, Article 1, 1--42.
\href{https://doi.org/10.5802/ojmo.49}
{doi:10.5802/ojmo.49}.

\end{thebibliography}
\end{document}
